% \newcommand{\lang}{french}  % Copy this in a french document
% \newcommand{\UseWhiteBackground}{1}  % Copy this to have a white document

\usepackage[utf8]{inputenc}
\usepackage{standalone}
\usepackage{etoolbox}
\ifdef{\lang}{}{\newcommand{\lang}{english}}  % If no \lang command is defined, default to english
\usepackage[\lang]{babel}
\usepackage{amssymb,amsmath,amsthm,epsfig,eepic,epic,verbatim,moreverb,fancybox}
\usepackage{stmaryrd}  % for \llbracket and \rrbracket
\usepackage{wasysym}
% \usepackage{bbold}  % NO! DON"T USE IT. It add mathbb for integers but makes it ugly for letters.
\usepackage{dsfont}
% \usepackage{tocstyle}  % To change the table of contents (TOC) style, and that's what we do next
% \usetocstyle{nopagecolumn}  % One of classic, KOMAlike, nopagecolumn, allwithdot, noonewithdot, standard
\usepackage{tikz}
\usepackage{tikz-cd}
\usetikzlibrary{calc,math,decorations.markings,decorations.pathmorphing,arrows,shadows,shapes.geometric,shapes.multipart,patterns,automata,backgrounds}
\usepackage{rotating}  % To have large figures in lanscape with \begin{sidewaysfigure} ... \end{sidewaysfigure}
\usepackage{pdflscape}  % To change one page to landscape
\usepackage{multicol}  % For the \begin{multicol}{N} environement
\usepackage{wrapfig}   % For the \begin{wrapfig}{r|l}{4cm} environement that behave like html float
\usepackage{subcaption}  % To use the subfigure environement in figure environements
\usepackage{multirow}  % Use \multirow{N}{*}{text} to have text in N rows in a table
\usepackage{mathtools}  % For the \xSOMETHINGarrow that we can write on top and change length
\usepackage{xspace}  % To add a space in commands when the following is not pounctuation
\usepackage{todonotes}  % To add todos in the margin of a document with \todo. Also \listoftodos is nice !
\usepackage{graphicx}  % many operations like rotating text etc...
\usepackage{ifthen}
\usepackage{xifthen}
\usepackage{float}  % To have the [H] option for figure so that they don't float. Use only if really needed.
\usepackage{pdfpages}
\usepackage{tablefootnote}  % To have footnotes in tables, use \tablefootnote{...}
\usepackage{enumerate}  % To be able to change the numbering of enumerate environements with [i)] or [a.]...
\usepackage{enumitem}  % To manage the spacing of items and margins
\usepackage{forest}  % For automatic tikz layout
\usepackage[framemethod=tikz]{mdframed}  % Frame around theorems
\usepackage{thmtools}  % Fixes the issue with \autoref and shared counters for thm env. Probably lots of other things also...
\usepackage{totcount}
\regtotcounter{figure}  % To count the number of figures, use \total{figure} to print it
\usepackage{listings}
% \usepackage{minted}
\usepackage{pbox}  % A variable width parbox
\usepackage{xcolor}

\definecolor{airforceblue}{rgb}{0.36, 0.54, 0.66}
\definecolor{awesome}{rgb}{1.0, 0.13, 0.32}
\definecolor{bondiblue}{rgb}{0.0, 0.58, 0.71}
\definecolor{atomictangerine}{rgb}{1.0, 0.6, 0.4}
\definecolor{ghostwhite}{rgb}{0.97, 0.97, 1.0}

% creamy background
\ifdef{\UseWhiteBackground}{  % White, default
    \definecolor{myBGcolor}{HTML}{FFFFFF}
    \definecolor{myTextcolor}{HTML}{000000}
    \definecolor{linkColor}{HTML}{044F67}

    % \colorlet{linkColor}{orange!60!blue}
    \colorlet{myRed}{awesome!10!myBGcolor}
    \colorlet{myGreen}{airforceblue!15!myBGcolor}
}{  % Cream
    \definecolor{myBGcolor}{HTML}{E6E0C6}
    \definecolor{myTextcolor}{HTML}{2F251C}
    \colorlet{linkColor}{orange!70!blue}

    \pagecolor{myBGcolor}
    \color{myTextcolor}

    \colorlet{myRed}{awesome!10!myBGcolor}
    \colorlet{myGreen}{green!10!myBGcolor}
}


\usepackage{hyperref}  % Always include it last, as it overrides a lot of commands.
\hypersetup{
    % bookmarks=true,
    % pdfpagemode=FullScreen,  % Only for presentations
    pdfauthor={Diego Dorn},
    pdftitle={ok},
    colorlinks=true,
    allcolors=linkColor
    % linkcolor=blue,
    % urlcolor={winered},
    % filecolor={winered},
    % citecolor={winered},
    % allcolors={orange},
    % linktoc=all,
}

% \usepackage{cleveref}  % Needs to be after hyperref

\addto\extrasenglish{%
    \renewcommand{\subsectionautorefname}{Section}  % Change name of subsection in autoref to be more uniform
}


% Tool to translate commands
% The first argument is the english text, the second is in french
\newcommand{\ifenglish}[2]{\ifthenelse{\equal{\lang}{english}}{#1}{#2}}

% Frames
\mdfdefinestyle{theoremstyle}{%
roundcorner=0pt, % round corners, makes it friendlier
leftmargin=1pt, rightmargin=1pt, % this helps with box warnings
innerbottommargin=12pt,
hidealllines=true, % even friendlier
align=center, %
}

% Theorem environements
\theoremstyle{definition}   % Not in italics
\ifenglish{
    \newtheorem{theorem}{Theorem}[section]
    \newtheorem*{theorem*}{Theorem}
    \newtheorem{theoremx}{Theorem}
    \newtheorem{definition}[theorem]{Definition}
    \newtheorem{notation}[theorem]{Notation}
    \newtheorem{corollary}[theorem]{Corollary}
    \newtheorem{proposition}[theorem]{Proposition}
    \newtheorem*{proposition*}{Proposition}
    \newtheorem{lemma}[theorem]{Lemma}
    \newtheorem*{lemma*}{Lemma}
    \newtheorem*{remark}{Remark}
    \newtheorem*{claim}{Claim}
    \newtheorem*{example}{Example}
    \newtheorem*{examples}{Examples}
    \newtheorem*{exercise}{Exercise}
}{
    \newtheorem{theorem}{Théorème}[section]
    \newtheorem*{theorem*}{Théorème}
    \newtheorem{theoremx}{Théorème}
    \newtheorem{definition}[theorem]{Définition}
    \newtheorem{notation}[theorem]{Notation}
    \newtheorem{corollary}[theorem]{Corrolaire}
    \newtheorem{proposition}[theorem]{Proposition}
    \newtheorem{lemma}[theorem]{Lemme}
    \newtheorem*{remark}{Remarque}
    \newtheorem*{claim}{Assertion}
    \newtheorem*{example}{Exemple}
    \newtheorem*{examples}{Exemples}
    \newtheorem*{exercise}{Exercise}
}
\renewcommand{\thetheoremx}{\Alph{thmx}} % "letter-numbered" theorems


\surroundwithmdframed[style=theoremstyle,backgroundcolor=myRed]{theorem}
\surroundwithmdframed[style=theoremstyle,backgroundcolor=myRed]{theorem*}
\surroundwithmdframed[style=theoremstyle,backgroundcolor=myGreen]{definition}

%%%%%%%%%%%%%%%%%%%
%   Parenthesis   %
%%%%%%%%%%%%%%%%%%%

\newcommand{\parenthesis}[1]{\left(#1\right)}
\newcommand{\bigparenthesis}[1]{\big(#1\big)}
\newcommand{\Bigparenthesis}[1]{\Big(#1\Big)}
\newcommand{\Par}[2][]{
    \ifthenelse{\isempty{#1}}{%
        \mathopen{}\left(#2\right)\mathclose{}%
    }{\ifthenelse{#1 = 0}{
        (#2)%
    }{\ifthenelse{#1 = 1}{%
        \big(#2\big)%
    }{\ifthenelse{#1 = 2}{%
        \Big(#2\Big)%
    }{\ifthenelse{#1 = 3}{%
        \large(#2\large)%
    }{\ifthenelse{#1 = 4}{%
        \Large(#2\Large)%
    }{%
        \mathopen{}\left(#2\right)\mathclose{}%
    }}}}}}}
\newcommand{\Bra}[2][]{
    \ifthenelse{\isempty{#1}}{%
        \mathopen{}\left[#2\right]\mathclose{}%
    }{\ifthenelse{#1 = 0}{
        [#2]%
    }{\ifthenelse{#1 = 1}{%
        \big[#2\big]%
    }{\ifthenelse{#1 = 2}{%
        \Big[#2\Big]%
    }{\ifthenelse{#1 = 3}{%
        \large[#2\large]%
    }{\ifthenelse{#1 = 4}{%
        \Large[#2\Large]%
    }{%
        \mathopen{}\left[#2\right]\mathclose{}%
    }}}}}}}


\newcommand{\brackets}[1]{\left[#1\right]}

%%%%%%%%%%%%%%%%%%%
%    Operators    %
%%%%%%%%%%%%%%%%%%%

\DeclareMathOperator{\id}{id}
\DeclareMathOperator{\re}{Re}
\DeclareMathOperator{\Aut}{Aut}
\DeclareMathOperator{\Hom}{Hom}
\DeclareMathOperator{\Char}{char}
\DeclareMathOperator{\im}{Im}
\DeclareMathOperator{\dom}{dom}
\DeclareMathOperator{\cof}{cof}
\DeclareMathOperator{\Ob}{Ob}
\DeclareMathOperator{\Mor}{Mor}
\DeclareMathOperator{\rk}{rk}

\newcommand{\concat}{
  \mathord{
    \mathchoice
    {\raisebox{1ex}{\scalebox{.7}{$\frown$}}}
    {\raisebox{1ex}{\scalebox{.7}{$\frown$}}}
    {\raisebox{.7ex}{\scalebox{.5}{$\frown$}}}
    {\raisebox{.7ex}{\scalebox{.5}{$\frown$}}}
  }
}
\newcommand{\restr}[1]{\mathbin\upharpoonright_{#1}}
\newcommand{\compl}[1]{#1^\mathsf{c}}
\newcommand{\abs}[1]{\left| #1 \right|}
\newcommand{\bigabs}[1]{\big\lvert#1\big\rvert}
\newcommand{\parens}[1]{\left(#1\right)}
\newcommand{\norm}[1]{\left\| #1 \right\|}

%%%%%%%%%%%%%%%%%%%
%  Constructions  %
%%%%%%%%%%%%%%%%%%%

% Arrows
\newcommand{\xto}[1]{\xrightarrow{#1}}

% Structures
\newcommand{\verteqsign}{\rotatebox{90}{$\,=$}}
\newcommand{\verteq}[2]{\underset{\scriptstyle\overset{\mkern4mu\verteqsign}{#2}}{#1}}
\newcommand{\vertimplies}[2]{\underset{\scriptstyle\overset{\uparrow}{#2}}{#1}}
\newcommand{\func}[5]{
    \begin{array}{rccl}
            #1 : & #2 & \longrightarrow & #3 \\
                 & #4 & \longmapsto     & #5
    \end{array}
}
\newcommand{\ntimes}[2]{\underbrace{#1, \dots, #1}_{#2 \ttimes}}
% Reminder: for combinations / binomial coefficients, use \binom{n}{k}

% Set definitions
\newcommand{\set}[1]{\left\{ #1 \right\}}
\newcommand{\bigset}[1]{\big\{ #1 \big\}}
\newcommand{\setst}[2]{\left\{ #1 \, \middle| \, #2 \right\}}
\newcommand{\bigsetst}[2]{\big\{ #1 \, \big| \, #2 \big\}}

\newcommand{\funcset}[2]{{}^{#1}{#2}}
\newcommand{\angles}[1]{\left\langle #1 \right\rangle}
\newcommand{\brz}{B_r(z_0)}
\newcommand\quotient[2]{
    \mathchoice
        {% \displaystyle
            \text{\raise1ex\hbox{$#1$}\Big/\lower1ex\hbox{$#2$}}%
        }
        {% \textstyle
            #1\,/\,#2
        }
        {% \scriptstyle
            #1\,/\,#2
        }
        {% \scriptscriptstyle
            #1\,/\,#2
        }
}
\newcommand{\interior}[1]{\mathring{#1}}

%%%%%%%%%%%%%%%
%   Letters   %
%%%%%%%%%%%%%%%


% Greek alphabe
\renewcommand{\phi}{\ensuremath\varphi\xspace}
\renewcommand{\epsilon}{\ensuremath\varepsilon\xspace}
\renewcommand{\a}{\ensuremath\alpha\xspace}
\renewcommand{\d}{\ensuremath{\partial}\xspace}
\newcommand{\D}{\ensuremath\Delta\xspace}
\newcommand{\e}{\ensuremath\epsilon\xspace}
\newcommand{\g}{\ensuremath\gamma\xspace}
\newcommand{\y}{\g}  % I mistake them every other time
\newcommand{\p}{\phi\xspace}
\newcommand{\s}{\sigma\xspace}
\newcommand{\w}{\ensuremath\omega\xspace}

% Mathbb alphabet

\newcommand{\IB}{\ensuremath{\mathbb B}\xspace}
\newcommand{\C}{\ensuremath{\mathbb C}\xspace}
\newcommand{\ID}{\ensuremath{\mathbb D}\xspace}
\newcommand{\IE}{\ensuremath{\mathbb E}\xspace}
\newcommand{\F}{\ensuremath{\mathbb F}\xspace}
\newcommand{\IG}{\ensuremath{\mathbb G}\xspace}
\newcommand{\N}{\ensuremath{\mathbb N}\xspace}
\newcommand{\R}{\ensuremath{\mathbb R}\xspace}
\newcommand{\IP}{\ensuremath{\mathbb P}\xspace}
\newcommand{\IS}{\ensuremath{\mathbb S}\xspace}
\newcommand{\IT}{\ensuremath{\mathbb T}\xspace}
\newcommand{\Q}{\ensuremath{\mathbb Q}\xspace}
\newcommand{\Z}{\ensuremath{\mathbb Z}\xspace}
\newcommand{\ONE}{\text{\usefont{U}{bbold}{m}{n}1}}

% Fancy fonts
\newcommand{\Aa}{\ensuremath{\mathcal A}\xspace}
\newcommand{\Bb}{\ensuremath{\mathcal B}\xspace}
\newcommand{\Cc}{\ensuremath{\mathcal C}\xspace}
\newcommand{\Dd}{\ensuremath{\mathcal D}\xspace}
\newcommand{\Ee}{\ensuremath{\mathcal E}\xspace}
\newcommand{\Ff}{\ensuremath{\mathcal F}\xspace}
\newcommand{\Gg}{\ensuremath{\mathcal G}\xspace}
\newcommand{\Hh}{\ensuremath{\mathcal H}\xspace}
\newcommand{\Ii}{\ensuremath{\mathcal I}\xspace}
\newcommand{\Jj}{\ensuremath{\mathcal J}\xspace}
\newcommand{\Kk}{\ensuremath{\mathcal K}\xspace}
\newcommand{\Ll}{\ensuremath{\mathcal L}\xspace}
\newcommand{\Mm}{\ensuremath{\mathcal M}\xspace}
\newcommand{\Nn}{\ensuremath{\mathcal N}\xspace}
\newcommand{\Oo}{\ensuremath{\mathcal O}\xspace}
\newcommand{\Pp}{\ensuremath{\mathcal P}\xspace}
\newcommand{\Qq}{\ensuremath{\mathcal Q}\xspace}
\newcommand{\Rr}{\ensuremath{\mathcal R}\xspace}
\newcommand{\Ss}{\ensuremath{\mathcal S}\xspace}
\newcommand{\Tt}{\ensuremath{\mathcal T}\xspace}
\newcommand{\Uu}{\ensuremath{\mathcal U}\xspace}
\newcommand{\Vv}{\ensuremath{\mathcal V}\xspace}
\newcommand{\Ww}{\ensuremath{\mathcal W}\xspace}
\newcommand{\Xx}{\ensuremath{\mathcal X}\xspace}
\newcommand{\Yy}{\ensuremath{\mathcal Y}\xspace}
\newcommand{\Zz}{\ensuremath{\mathcal Z}\xspace}
\renewcommand{\O}[1]{\mathcal{O}\left(#1\right)}

% Others
\renewcommand{\emptyset}{\varnothing}


%%%%%%%%%%%%%%%%%%
%   New symbols  %
%%%%%%%%%%%%%%%%%%

\newcommand{\iddots}{{\cdot^{\cdot^\cdot}}}  % Three dots on the antidiagonal
\newcommand{\existsinf}{\exists^\infty}
\newcommand{\coloneq}{\mathrel{\resizebox{\widthof{$\mathord{=}$}}{\height}{ $\!\!\resizebox{1.2\width}{0.8\height}{\raisebox{0.23ex}{$\mathop{:}$}}\!\!=\!\!$ }}}  % Fancier :=

%%%%%%%%%%%%%%%%%%
%   Shorthands   %
%%%%%%%%%%%%%%%%%%

% Math text
\newcommand{\open}{\text{ open }}
\newcommand{\st}{\text{ \ifenglish{s.t}{t.q}. }}
\newcommand{\andd}{\text{ \ifenglish{and}{et} }}
\newcommand{\andquad}{\quad\andd\quad}
\newcommand{\et}{\text{ et }}
\newcommand{\ttimes}{\text{ times}}
\newcommand{\tif}{\text{if }}
\newcommand{\otherwise}{\text{otherwise}}

% Symbols
\newcommand{\acts}{\circlearrowright}
\newcommand{\sub}{\subset}
\newcommand{\subeq}{\subseteq}
\renewcommand{\iff}{\ensuremath{\Longleftrightarrow}}
\newcommand{\iffquad}{\quad\iff\quad}
\newcommand{\del}{\partial}
\newcommand{\x}{\times}
\renewcommand{\emph}{\textbf}
\newcommand{\Forces}{\vdash}

%%%%%%%%%%%%%%%%%%%%
% Subject specific %
%%%%%%%%%%%%%%%%%%%%

% Maybe those should go inside their respesective files

% Graph theory
\DeclareMathOperator*{\ex}{ex}
\newcommand{\Turan}{\mathcal{T}}

% Analytic number theory
\newcommand{\sumdn}{\sum_{d|n}}
\newcommand{\Li}{\mathrm{Li}}

% Smooth manifolds
\newcommand{\chart}[1]{(\phi_{#1}, U_{#1})}
\newcommand{\der}[1]{\frac{\del}{\del #1}}

% Galois Theory
\newcommand{\clot}[1]{\overline{#1}}
\newcommand{\Gal}{\mathrm{Gal}}
% \DeclareMathOperator*{\Gal}{Gal}
\newcommand{\Kxx}{K[x_1, \dots, x_n]}

% Algebraic topology
\newcommand{\HCW}{H^{\mathrm{CW}}}

% Complexity
\newcommand{\coP}{\mathsf{coP}}
\newcommand{\Ptime}{\mathsf{P}}
\newcommand{\PHierachy}{\mathsf{PH}}
\newcommand{\NP}{\mathsf{NP}}
\newcommand{\coNP}{\mathsf{coNP}}
\newcommand{\PSPACE}{\mathsf{PSPACE}}
\newcommand{\NPSPACE}{\mathsf{NPSPACE}}

% Analysis
\newcommand{\dd}{\mathrm{d}}

% Probability

% Neural networks
\newcommand{\MSE}{\mathrm{MSE}}
\newcommand{\KLdiv}[3][]{D_{\mathrm{KL}}\Par[#1]{#2 \mid\mid #3}}
\DeclareMathOperator{\Sign}{Sign}
% A Devil emoji in exponent of the argument
\newcommand{\Adversarial}[1]{#1^{\raisebox{-0.2ex}{\includegraphics[height=2ex]{images/devil.png}}}}

% Logic
\newcommand{\Symbol}[1]{\underline{#1}}

% Descriptive Set Theory
\newcommand{\baire}{\ensuremath{\w^\w}\xspace}
\newcommand{\cantor}{\ensuremath{2^\w}\xspace}
\newcommand{\W}{\ensuremath{\w_1}\xspace}  % This is meant to be used only for the first uncountable ordinal
% Classes
\newcommand{\bS}[1]{\ensuremath{\mathbf{\Sigma}^0_{#1}}\xspace}
\newcommand{\bP}[1]{\ensuremath{\mathbf{\Pi}^0_{#1}}\xspace}
\newcommand{\bD}[1]{\ensuremath{\mathbf{\Delta}^0_{#1}}\xspace}
\newcommand{\borelsigma}[1]{\ensuremath{\mathbf{\Sigma}^0_{#1}}\xspace}
\newcommand{\borelpi}[1]{\ensuremath{\mathbf{\Pi}^0_{#1}}\xspace}
\newcommand{\boreldelta}[1]{\ensuremath{\mathbf{\Delta}^0_{#1}}\xspace}
\newcommand{\borelGamma}{\ensuremath{\mathbf{\Gamma}}\xspace}
\newcommand{\LambdaClass}[2]{\ensuremath{\mathbf{\Lambda}_{#1, #2}}\xspace}
\newcommand{\dualclass}[1]{\check{#1}}
\newcommand{\nsd}{non-self-dual\xspace}
\newcommand{\projectivesigma}[1]{\ensuremath{\mathbf{\Sigma}^1_{#1}}\xspace}
\newcommand{\projectivepi}[1]{\ensuremath{\mathbf{\Pi}^1_{#1}}\xspace}
\newcommand{\projectivedelta}[1]{\ensuremath{\mathbf{\Delta}^1_{#1}}\xspace}
% Sets of sequences
\newcommand{\seqlt}[2]{{#2}^{<#1}}
\newcommand{\finiteseq}[1]{\seqlt{\w}{#1}}
% \newcommand{\infseq}[1]{\funcset{\omega}{#1}}
% \newcommand{\allseq}[1]{\funcset{\leq\omega}{#1}}
\newcommand{\infseq}[1]{{#1}^\w}
\newcommand{\allseq}[1]{{#1}^{\leq\omega}}
\newcommand{\finitebaire}{\ensuremath{\finiteseq{\w}}\xspace}
\newcommand{\finitecantor}{\ensuremath{\finiteseq{2}}\xspace}
\newcommand{\length}[1]{\mathrm{lh}(#1)}
\newcommand{\emptyseq}{\langle\rangle}
% Operators
\renewcommand{\succ}[2][]{\ifthenelse{\isempty{#1}}{\mathrm{succ}(#2)}{\mathrm{succ}_{#1}(#2)}}
\DeclareMathOperator{\odd}{odd}
\DeclareMathOperator{\even}{even}
\DeclareMathOperator{\last}{last}
\newcommand{\Countf}[1]{\#_{#1}}
\newcommand{\difference}[1]{D_{#1}}
\newcommand{\suffixes}[2]{{#1}_{(#2)}}
\newcommand{\wadd}{\oplus}
\newcommand{\wCountableSup}{\mathrm{sup^\w}}
\newcommand{\wmult}{\mathop{\otimes}}
\newcommand{\bigeraser}{\ensuremath{\mathord{\Uparrow}}\xspace}
\renewcommand{\interleave}{\mathbin{\S}}
\newcommand{\Init}[1]{\mathrm{Init}_{#1}}
% Players
\newcommand{\I}{\texttt{I}\xspace}
\newcommand{\II}{\texttt{II}\xspace}
\newcommand{\playerA}{Player \I}
\newcommand{\playerB}{Player \II}

% \renewcommand{\I}{\texttt{M}\xspace}
% \renewcommand{\II}{\texttt{B}\xspace}
% \renewcommand{\playerA}{Michael\xspace}
% \renewcommand{\playerB}{Bob\xspace}
% Wadge
\newcommand{\leW}{\le_\textsc{w}}
\newcommand{\geW}{\ge_\textsc{w}}
\newcommand{\gtW}{>_\textsc{w}}
\newcommand{\ltW}{<_\textsc{w}}
\newcommand{\equivW}{\equiv_\textsc{w}}
\newcommand{\Wdegree}[1]{[#1]_{\equivW}}
% Axioms / Theories
\newcommand{\ZF}{\textbf{ZF}\xspace}
\newcommand{\ZFC}{\textbf{ZFC}\xspace}
\newcommand{\AD}{\textbf{AD}\xspace}
\renewcommand{\AC}{\textbf{AC}\xspace}
\newcommand{\DC}{\textbf{DC}\xspace}
% Games
\newcommand{\GaleStewart}[1]{\ensuremath{G(#1)}\xspace}
\newcommand{\BanachMazur}[1]{\ensuremath{BM(#1)}\xspace}
\newcommand{\Wadge}[2]{\ensuremath{G(#1, #2)}\xspace}
% Automatons
\newcommand{\automaton}{automatic set\xspace}
\newcommand{\automata}{automatic sets\xspace}
\newcommand{\Automata}{Automatic sets\xspace}
\newcommand{\Language}{\mathcal{L}}
\DeclareMathOperator{\Inf}{Inf}
\newcommand{\AcceptanceCondition}{\mathrm{Acc}}
\newcommand{\automatonMinus}{\scalebox{0.5}{\begin{tikzpicture}[automata, baseline=-2mm]
    \node[minus] {};
\end{tikzpicture}}\xspace}
\newcommand{\automatonPlus}{\scalebox{0.5}{\begin{tikzpicture}[automata, baseline=-2mm]
    \node[plus] {};
\end{tikzpicture}}\xspace}
\newcommand{\automatonState}[1]{\scalebox{0.5}{\begin{tikzpicture}[automata, baseline=-2mm, minimum size=20]
    \node[state] {\Large \(#1\)};
\end{tikzpicture}}\xspace}
\newcommand{\sumAutomaton}[2]{\begin{tikzpicture}[automata]
    \path node[initial, state] (root) {$#2$}
        +(2, 1) node[state] (A) {$#1$}
        +(2, -1) node[state] (Ac) {$\compl{#1}$};
    \draw (root) edge (A) edge (Ac);
\end{tikzpicture}}


%%%%%%%%%%%%%%%%%
%     Tikz      %
%%%%%%%%%%%%%%%%%

\tikzcdset{row sep/normal=1cm, column sep/normal=1cm}  % So tikz-cd diagram are more squared

% Draws irregular circles.
\newcommand{\irregularcircle}[2]{  % radius, irregularity
  let \n1 = {(#1)+rand*(#2)} in
  +(0:\n1)
  \foreach \a in {10,20,...,350}{
    let \n1 = {(#1)+rand*(#2)} in
    -- +(\a:\n1)
  } -- cycle
}

\tikzset{
triangle/.style={ % To create triangles in trees
  draw=myRed!90!black,
  text=black,
  fill=myRed,
  shape border uses incircle,
  isosceles triangle,
  shape border rotate=90,
}}

% Used to make arrows on paths. Use like this: \draw[marr=\Singlearrow]
\tikzset{marr/.style={
  decoration={
    markings,
    mark=at position 0.5 with {#1}
    },
  postaction={decorate}
}}
\def\Singlearrow{{\arrow[scale=1.5,xshift={0.5pt+2.25\pgflinewidth}]{>}}}
\def\Doublearrow{{\arrow[scale=1.5,xshift=1.35pt +2.47\pgflinewidth]{>>}}}
\def\Triplearrow{{\arrow[scale=1.5,xshift=1.75pt +2.47\pgflinewidth]{>>>}}}
\newcommand{\CWarrow}[3]{ [->] ($1/3*(#1) + 1/3*(#2) + 1/3*(#3)$) ++(220:4pt) arc (220:-40:4pt)}
\newcommand{\CCWarrow}[3]{ [->] ($1/3*(#1) + 1/3*(#2) + 1/3*(#3)$) ++(-40:4pt) arc (-40:220:4pt)}

% For automatons
\tikzset{
    initial text={},
    do path picture/.style={%
        path picture={%
        \pgfpointdiff{\pgfpointanchor{path picture bounding box}{south west}}%
            {\pgfpointanchor{path picture bounding box}{north east}}%
        \pgfgetlastxy\x\y%
        \tikzset{x=\x/2,y=\y/2}%
        #1
        }
    },
    automata/.style={
        node distance=60pt,
        minimum size=20,
        every edge/.style={
            draw,
            ->,
            auto,
        }
    },
    plus/.style={
        draw,
        circle,
        do path picture={
            \draw [line cap=round, very thick] (-1/2,0) -- (1/2,0) (0,-1/2) -- (0,1/2);
        }
    },
    minus/.style={
        draw,
        circle,
        do path picture={
            \draw [line cap=round, very thick] (-1/2,0) -- (1/2,0);
        }
    },
    % Removes extra space on initial nodes
    % From: https://tex.stackexchange.com/questions/111554/how-to-avoid-superfluous-space-in-the-initial-state-of-a-tikz-automaton#comment245757_111558
    every initial by arrow/.append style={anchor/.append style={shape=coordinate}},
    state/.append style={fill=white},
}

% To put at the bottom right corner of a square that commutes.
\newcommand{\commutes}{\arrow[ul, phantom, "\scalebox{1.5}{$\circlearrowleft$}"]}

% To link two lines in a long exact sequence.
% First argument is where the line should connect (ex: dll)
% Second argument is the label on the line.
\newcommand{\connecting}[2]{
    \arrow[#1,phantom, ""{coordinate, name=Z}]
    \arrow[#1, "#2", rounded corners,
        to path={
            -- ([xshift=2ex]\tikztostart.east)
            |- (Z)[near start]\tikztonodes
            -| ([xshift=-2ex]\tikztotarget.west)
            -- (\tikztotarget)
        }]
}
\newcommand{\dotsto}{\mathllap{\cdots\longrightarrow\ }}  % ... → indicate the start of a long sequence without taking space
\newcommand{\todots}{\mathrlap{\ \longrightarrow\cdots}}  % → ... indicate the end of a long sequence without taking space

\newenvironment{longsequence}{
    \begin{center}
        \begin{tikzcd}[row sep=0.5cm]
}{
        \end{tikzcd}
    \end{center}
}


\newcommand{\quickfig}[2]{
    \begin{wrapfigure}{r}{40mm}
        \begin{center}
            #2
        \end{center}
        \caption{#1}
    \end{wrapfigure}
}
